\Article{5}
Die Formel, mit der wir im vorigen Artikel das Produkt beliebig vieler Funktionen ausgedrückt haben, ist allgemein und setzt keinen Zusammenhang zwischen den Wurzeln $r,r',r'',r'''$ usw.\ oder zwischen $R,R',R'',R'''$ usw.\ voraus. Daraus lässt sich leicht ableiten: Dürfen die Wurzeln $r',r'',r'''$ usw.\ als Potenzen von $r$ und die Wurzeln $R',R'',R'''$ usw.\ als Potenzen von $R$ betrachtet werden, so hat jedes Glied des Produkts die Form $R^M[r^m,R^\lambda]$, wobei der Exponent $\lambda$ für alle Glieder derselbe ist, nämlich $R^\lambda=RR'R''R'''$ usw. Nach dem, was wir in Artikel~3 bemerkt haben, lässt sich daher ein solches Produkt auf die folgende Form bringen:
\[
\begin{aligned}
&A[1,R^\lambda]+B[r,R^\lambda]+B'[r^g,R^\lambda]+B''[r^{g^2},R^\lambda]+B'''[r^{g^3},R^\lambda]+\text{usw.}\\
&\hspace{5em}{}+B^{(\alpha-1)}[r^{g^{\alpha-1}},R^\lambda]
\end{aligned}
\]
Dabei haben die einzelnen Koeffizienten $A,B,B',B'',B'''$ usw.\ die Form
\[
h+h'R+h''R^2+h'''R^3+\text{usw.}+h^{(\beta-1)}R^{\beta-1}
\]
und $h,h',h'',h'''$ usw.\ bezeichnen bestimmte ganze Zahlen.

\SourceLeaf{257}
Der einfachste Fall ist der, in dem $r=r'=r''=r'''$ usw.\ und auch $R=R'=R''=R'''$ usw.\ gesetzt wird; dann geht unser Produkt in die Potenz $[r,R]^\lambda$ über, die sich daher stets auf die oben angegebene Form bringen lässt.

\Article{6}
Setzen wir also $\lambda=\beta$, so erhält die Potenz $[r,R]^\beta$ die Form
\[
\begin{aligned}
&A[1,1]+B[r,1]+B'[r^g,1]+\text{usw.}+B^{\alpha-1)}[r^{g^{\alpha-1}},1]\\
&{}=\beta\gamma A+B(\beta\gamma,1)+B'(\beta\gamma,g)+B''(\beta\gamma,g^2)+\text{usw.}\\
&\qquad{}+B^{(\alpha-1)}(\beta\gamma,g^{\alpha-1})=\theta'
\end{aligned}
\]

Werden also nicht nur der Wert der Wurzel $R$ (und damit auch die Werte der Koeffizienten $A,B,B'$ usw.), sondern auch die Werte der einzelnen Summen von $\beta\gamma$ Gliedern $(\beta\gamma,1),(\beta\gamma,g)$ usw.\ als bekannt vorausgesetzt, so ist der Wert von $\theta'$ unmittelbar bekannt, und daraus lässt sich $[r,R]$ durch die Formel ${}^{\beta}\!\surd\theta'$ ermitteln. Dieser Ausdruck lässt $\beta$ verschiedene Werte zu; daher könnte zweifelhaft erscheinen, welchen man nehmen soll. Man zeigt jedoch leicht, dass diese Wahl völlig beliebig ist, sofern $R$ eine \emph{eigentliche} Wurzel der Gleichung $x^\beta-1=0$ ist. In diesem Fall sind nämlich jene $\beta$ Werte des Wurzelausdrucks ${}^{\beta}\!\surd\theta'$ offenbar
\[
[r,R],\quad[r^{g^\alpha},R],\quad[r^{g^{2\alpha}},R]\dotsfour[r^{g^{\alpha\beta-\alpha}},R]
\]
denn die Potenzen dieser Funktionen mit dem Exponenten $\beta$ sind nach Artikel~3 einander gleich, während die Funktionen selbst zu den $\beta$ verschiedenen Wurzeln der Gleichung $x^\beta-1=0$ proportional sind. Solange aber nur die Summen von $\beta\gamma$ Gliedern $(\beta\gamma,1),(\beta\gamma,g)$ usw.\ bekannt sind, ist die Wurzel $r$ selbst nur insofern bestimmt, als sie in der Menge $(\beta\gamma,1)$ enthalten sein muss; es bleibt beliebig, welche Wurzel aus dieser Menge wir für $r$ nehmen. Diese Wurzeln sind $r,r^{g^\alpha},r^{g^{2\alpha}}$ usw., und daher können wir auch von den Funktionen $[r,R],[r^{g^\alpha},R],[r^{g^{2\alpha}},R]$ usw.\ eine beliebige für $[r,R]$ nehmen.

Diese Schlüsse würden nicht gelten, wenn $R$ keine eigentliche Wurzel der Gleichung $x^\beta-1=0$ wäre. Nimmt man nämlich an, dass $R$ eine eigentliche Wurzel der Gleichung $x^{\beta'}-1=0$ und $\beta'$ ein Teiler von $\beta$ ist, so ergibt sich leicht
\[
[r,R]=[r^{g^{\alpha\beta'}},R],\quad[r^{g^\alpha},R]=[r^{g^{\alpha\beta'+\alpha}},R]\text{ usw.}
\]
Daher kommen in der Menge der $\beta$ Funktionen $[r,R],[r^{g^\alpha},R]\ldots[r^{g^{\alpha\beta-\alpha}},R]$ nur $\beta'$ verschiedene vor; folglich sind auch von den Werten des Ausdrucks ${}^{\beta}\!\surd\theta'$ nicht mehr als $\beta'$ zulässig, die übrigen $\beta-\beta'$ dagegen unzulässig. Man sieht aber leicht, dass man in diesem Fall nicht bis zur Potenz mit dem Exponenten $\beta$ der Funktion $[r,R]$
\SourceLeaf{258}
aufsteigen muss, sondern dass sich bereits die Potenz $[r,R]^{\beta'}$ auf unsere Form
\[
\beta\gamma A+B(\beta\gamma,1)+B'(\beta\gamma,g)+B''(\beta\gamma,g^2)\text{ usw.}
\]
bringen lässt. Wir erhalten somit $[r,R]$ durch einen Ausdruck wie ${}^{\beta'}\!\surd\theta$, und es ist gleichgültig, welchen Wert dieses Ausdrucks wir nehmen.

\Article{7}
Ebenso wie $[r,R]$ lassen sich auch die Funktionen $[r,R^2]$, $[r,R^3]$ usw.\ oder allgemein $[r,R^k]$ bestimmen. Nimmt man nämlich an, dass in $\theta'$ anstelle von $R$ die Potenzen $R^2,R^3$ usw.\ $R^k$ eingesetzt werden und dadurch die Funktionen $\theta'',\theta'''$ usw.\ $\theta^{(k)}$ entstehen, so ist offenbar $[r,R^2]^\beta=\theta''$, $[r,R^3]^\beta=\theta'''$ usw.\ und allgemein $[r,R^k]^\beta=\theta^{(k)}$; daher können auch diese Funktionen durch Wurzelausdrücke dargestellt werden, $[r,R^2]={}^{\beta}\!\surd\theta''$ usw. Es ist jedoch nicht zweckmäßig, diese Wurzelausdrücke zu verwenden, wenn eine Größe durch eine Funktion von $[r,R]$, $[r,R^2]$ usw.\ ausgedrückt werden soll. Da nämlich ihre einzelnen Werte nicht vollständig bestimmt sind, bliebe zweifelhaft, welche man miteinander kombinieren dürfte; offenbar ist dies aber keineswegs beliebig. Man sieht leicht, dass, sobald für $[r,R]$ ein bestimmter Wert genommen wird, auch sämtliche $[r,R^2]$, $[r,R^3]$ usw.\ vollständig bestimmte Werte erhalten müssen, die jedoch durch die Wurzelausdrücke nicht angegeben werden. Diese sind daher beiseitezusetzen und andere Ausdrücke zu suchen, mit deren Hilfe $[r,R^2]$, $[r,R^3]$ usw.\ \emph{rational} durch $[r,R]$ und bekannte Größen dargestellt werden. Dies erreichen wir leicht auf folgende Weise.

Nach dem Satz aus Artikel~4 und dem, was wir in Artikel~5 gezeigt haben, lässt sich auch das Produkt $[r,R^k]\times[r,R]^{\beta-k}$ auf eine Form wie
\[
\begin{aligned}
&\beta\gamma A+B(\beta\gamma,1)+B'(\beta\gamma,g)+B''(\beta\gamma,g^2)+\text{usw.}\\
&\qquad{}+B^{(\alpha-1)}(\beta\gamma,g^{\alpha-1})
\end{aligned}
\]
bringen, wobei $A,B,B',B''$ usw.\ rationale Funktionen von $R$ sind. Setzen wir also die Produkte
\[
\begin{aligned}
\relax[r,R^2]\times[r,R]^{\beta-2}&=\vartheta''\\
\relax[r,R^3]\times[r,R]^{\beta-3}&=\vartheta'''\\
\relax[r,R^4]\times[r,R]^{\beta-4}&=\beta''''\\
&\text{usw.}
\end{aligned}
\]
so sind auch $\vartheta'',\vartheta''',\vartheta''''$ usw.\ rational bestimmbare Größen, und es gilt
\SourceLeaf{259}
\[
\begin{aligned}
\relax[r,R^2]&=\frac{\vartheta''}{\theta'}[r,R]^2\\
\relax[r,R^3]&=\frac{\vartheta'''}{\theta'}[r,R]^3\\
\relax[r,R^4]&=\frac{\vartheta''''}{\theta'}[r,R]^4\\
&\text{usw.}
\end{aligned}
\]

\Needspace{10\baselineskip}Diese Ausdrücke stellen somit die Werte der Funktionen $[r,R^2]$, $[r,R^3]$ usw.\ rational dar, sofern nicht $[r.R]=0$ ist; gälte diese Gleichheit, so würden sie unbestimmt. Wir können aber streng beweisen, dass $[r,R]=0$ niemals eintreten kann, sofern $r$ eine von $1$ verschiedene Wurzel bezeichnet, auch wenn wir uns die Darlegung dieses Beweises für eine andere Gelegenheit vorbehalten müssen, um hier nicht zu ausführlich zu werden.

\Article{8}
Was wir in den vorigen Artikeln dargelegt haben, ist nützlich, wenn wir von Perioden mit $\beta\gamma$ Gliedern zu Perioden mit $\gamma$ Gliedern absteigen wollen. Man sieht nämlich leicht, dass, wenn $R$ eine eigentliche Wurzel bezeichnet, gilt:
\[
\begin{aligned}
\beta(\gamma,1)&=(\beta\gamma,1)+[r,R]+[r,R^2]+[r,R^3]+\text{usw.}+[r,R^{\beta-1}]\\
\beta(\gamma,g^\alpha)&=(\beta\gamma,1)+R^{\beta-1}[r,R]+R^{\beta-2}[r,R^2]\\
&\qquad{}+R^{\beta-3}[r,R^3]+\text{usw.}+R[r,R^{\beta-1}]\\
\beta(\gamma,g^{2\alpha})&=(\beta\gamma,1)+R^{2\beta-2}[r,R]+R^{2\beta-4}[r,R^2]\\
&\qquad{}+R^{2\beta-6}[r,R^3]+\text{usw.}+R^2[r,R^{\beta-1}]\\
&\text{usw.}
\end{aligned}
\]

Nähme man hier für die einzelnen $[r,R]$, $[r,R^2]$ usw.\ die Wurzelausdrücke ${}^{\beta}\!\surd\theta'$, ${}^{\beta}\!\surd\theta''$ usw., so wäre der Wert jeder Reihe unter $\beta^{\beta-1}$ Werten zweifelhaft. Nimmt man dagegen für $[r,R^2]$ usw.\ die rationalen Ausdrücke, so ist er keiner anderen Mehrdeutigkeit unterworfen als derjenigen, die der Natur der Sache nach unvermeidlich ist. Diese Bemerkung scheint der Aufmerksamkeit des berühmten \textsc{Lagrange} entgangen zu sein, der glaubte, unsere in den Disquis.\ arithm.\ Artikel~360 angegebene Methode \emph{vereinfacht} zu haben, indem er jene Wurzelausdrücke anstelle der rationalen Ausdrücke einsetzte; dabei hatten wir dort die Wurzelausdrücke mit Bedacht übergangen und nur rationale Ausdrücke angegeben (Traité de la résolution numérique des équations; Ausgabe $2$, S.\ 311).

Ferner muss hier kaum darauf hingewiesen werden, dass, sobald die Werte der Perioden $(\gamma,1)$, $(\gamma,g^\alpha)$ usw.\ oder auch nur einer von ihnen ermittelt sind, die Werte aller übrigen Perioden mit $\gamma$ Gliedern rational daraus abgeleitet werden können. Der Abstieg von Perioden mit $\beta\gamma$ Gliedern zu Perioden mit $\gamma$ Gliedern erfordert somit die Lösung der Gleichungen $x^\beta=1$, $x^\beta=\theta'$; die übrigen Operationen werden rational ausgeführt.
